%==============================================================================
% Chapitre 11 - Vitesse du son et nombre de Mach
%==============================================================================

\section{Introduction}

La vitesse du son joue un rôle fondamental en balistique.

De nombreux phénomènes aérodynamiques dépendent non seulement de la vitesse
du projectile, mais également de sa vitesse par rapport au son.

Cette comparaison est réalisée à l'aide du nombre de Mach.

La compréhension de ces notions est essentielle pour étudier :

\begin{itemize}
\item les projectiles supersoniques ;
\item les projectiles subsoniques ;
\item les ondes de choc ;
\item la traînée aérodynamique ;
\item les performances à longue distance.
\end{itemize}

\begin{aretenir}

Ce n'est pas la vitesse absolue du projectile qui importe, mais sa vitesse
par rapport à la vitesse du son dans l'air qui l'entoure.

\end{aretenir}

\section{Qu'est-ce que le son ?}

Le son est une onde mécanique.

Il se propage grâce aux vibrations des molécules du milieu traversé.

Contrairement à la lumière, le son ne peut pas se propager dans le vide.

Sa vitesse dépend des propriétés du milieu et notamment :

\begin{itemize}
\item de sa composition ;
\item de sa température ;
\item de sa pression ;
\item de sa densité.
\end{itemize}

\section{Vitesse du son dans l'air}

Dans l'air sec au niveau de la mer, la vitesse du son est d'environ :

\begin{equation}
c \approx 343~\mathrm{m/s}
\end{equation}

pour une température de 20°C.

Cette valeur varie avec la température.

\begin{rigueur}

Une approximation couramment utilisée est :

\begin{equation}
c \approx 331 + 0,6,T
\label{eq:vitesse_son}
\end{equation}

où :

\begin{itemize}
\item $c$ est la vitesse du son en m/s ;
\item $T$ est la température en °C.
\end{itemize}

\end{rigueur}

\section{Définition du nombre de Mach}

Le nombre de Mach est défini comme le rapport entre la vitesse d'un objet et
la vitesse du son.

\begin{equation}
M = \frac{v}{c}
\label{eq:mach}
\end{equation}

où :

\begin{itemize}
\item $M$ est le nombre de Mach ;
\item $v$ est la vitesse de l'objet ;
\item $c$ est la vitesse du son.
\end{itemize}

Le nombre de Mach est sans dimension.

\section{Interprétation du nombre de Mach}

Le nombre de Mach indique combien de fois un objet se déplace plus vite que
le son.

\begin{exemple}

Pour un projectile se déplaçant à :

\begin{equation}
v = 850~\mathrm{m/s}
\end{equation}

dans un air où :

\begin{equation}
c = 343~\mathrm{m/s}
\end{equation}

on obtient :

\begin{equation}
M \approx 2,48
\end{equation}

Le projectile vole donc à environ Mach 2,5.

\end{exemple}

\section{Régime subsonique}

Un objet est dit subsonique lorsque :

\begin{equation}
M < 1
\end{equation}

Sa vitesse est alors inférieure à celle du son.

Les perturbations aérodynamiques se propagent devant lui.

Les écoulements sont généralement relativement stables.

\section{Régime supersonique}

Un objet est dit supersonique lorsque :

\begin{equation}
M > 1
\end{equation}

Il se déplace alors plus rapidement que les perturbations qu'il génère.

Des ondes de choc apparaissent.

La plupart des munitions militaires modernes quittent le canon dans ce
régime.

\begin{exemple}

Une munition de 7,62~mm OTAN tirée à 850~m/s est initialement supersonique.

\end{exemple}

\section{Régime transsonique}

Le régime transsonique correspond à la zone située autour de Mach 1.

Dans ce manuel, nous considérerons généralement :

\begin{itemize}
\item subsonique : $M < 0,8$ ;
\item transsonique : $0,8 \leq M \leq 1,2$ ;
\item supersonique : $M > 1,2$.
\end{itemize}

Ces limites sont indicatives et peuvent varier selon les auteurs.

\begin{aretenir}

La zone transsonique est souvent la plus complexe à modéliser.

\end{aretenir}

\section{Pourquoi la zone transsonique est-elle importante ?}

Lorsqu'un projectile ralentit et approche de Mach 1 :

\begin{itemize}
\item les écoulements changent ;
\item les ondes de choc évoluent ;
\item les coefficients aérodynamiques peuvent varier rapidement ;
\item la stabilité peut être affectée.
\end{itemize}

Cette zone constitue souvent l'une des principales difficultés des modèles
balistiques.

\section{Ondes de choc}

Lorsqu'un projectile est supersonique, il génère des ondes de choc.

Ces ondes forment une surface conique appelée cône de Mach.

Le « bang » caractéristique associé au passage d'une balle supersonique est
directement lié à ces ondes de choc.

\begin{exemple}

Le bruit perçu lors du passage d'un projectile peut être entendu avant le
bruit de départ du tir.

Cette situation peut sembler paradoxale mais résulte simplement de la
propagation différente du projectile et du son.

\end{exemple}

\section{Influence sur la traînée}

Le coefficient de traînée d'un projectile varie fortement selon son nombre
de Mach.

Les zones les plus sensibles sont généralement :

\begin{itemize}
\item l'approche de Mach 1 ;
\item la zone transsonique ;
\item certaines vitesses supersoniques.
\end{itemize}

Cette dépendance explique pourquoi les modèles balistiques modernes utilisent
souvent des tables aérodynamiques complexes.

\section{Application aux essais}

Lors d'une campagne d'essais, il peut être utile de connaître :

\begin{itemize}
\item la vitesse initiale ;
\item la vitesse résiduelle ;
\item le nombre de Mach correspondant.
\end{itemize}

Ces informations permettent notamment d'identifier le régime de vol du
projectile.

\begin{simulation}

Les modèles balistiques avancés utilisent fréquemment le nombre de Mach pour
déterminer :

\begin{itemize}
\item le coefficient de traînée ;
\item certains coefficients aérodynamiques ;
\item les changements de régime de vol.
\end{itemize}

Le nombre de Mach est souvent l'une des variables principales des modèles
modernes.

\end{simulation}

\section{À retenir}

\begin{aretenir}

\begin{itemize}
\item Le nombre de Mach compare la vitesse d'un objet à la vitesse du son.
\item Un projectile peut être subsonique, transsonique ou supersonique.
\item La vitesse du son dépend notamment de la température.
\item Les ondes de choc apparaissent en régime supersonique.
\item La zone transsonique est particulièrement importante en balistique.
\end{itemize}

\end{aretenir}

